\section{A conversion workflow for existing proof packets}

The repository's packets differ greatly in maturity.  The following workflow
turns each into a comparable, auditable artifact.

\subsection{Step 1: state the constructive theorem}

Replace a classical slogan by an explicit input/output contract.  State:
the numerical input code, well-formedness test, request, output relation,
error metric, and every extra evidence item.  Separate unconditional
approximate output from gap-conditioned exact output.

\subsection{Step 2: classify every dependency}

Each dependency receives one of five labels:
\[
\begin{array}{ll}
\mathsf{PR}&\text{fixed primitive-recursive computation},\\
\mathsf{LOCAL}&\text{proved finite arithmetic contract},\\
\mathsf{CERT}&\text{supplied finite certificate with verifier},\\
\mathsf{PRODUCER}&\text{theorem that constructs such a certificate},\\
\mathsf{BARRIER}&\text{interface ruled out or known to require stronger data}.
\end{array}
\]
This prevents a missing producer from being mistaken for a coding problem.

\subsection{Step 3: build the contract graph}

Make each substantive lemma a node.  Edges record both mathematical
dependence and precision demand.  Common nodes---rational arithmetic,
polynomial evaluation, finite sums, bounded search, interval refinement,
finite nets---should point to shared libraries rather than being copied.

\subsection{Step 4: totalize operations}

Every operation must be total on all number codes.  Division, normalization,
chart inversion, and array access return a harmless default outside their
domain.  A Boolean guard is carried in the invariant and proves that the
default branch is never used in a valid run.

\subsection{Step 5: compute reverse demands}

Starting from the theorem's requested output precision, compute every
upstream precision, grid scale, cutoff, and iteration count.  Record these
functions explicitly.  Phrases such as ``take \(N\) large enough'' are not
accepted as a contract.

\subsection{Step 6: isolate discontinuities}

For every exact branch ask:
\begin{enumerate}[label=(\alph*)]
\item Can the output be quotient-invariant or relational?
\item Can overlapping guards make the approximate theorem uniform?
\item Which positive gap would make the exact branch decidable?
\item Is the requested selector contradicted by a known computability or
      constructive-logic reduction?
\end{enumerate}

\subsection{Step 7: define the finite trace invariant}

The invariant must mention only:
loop counters, coded states, Boolean verifiers, bounds, and finite folds.
If it contains ``for every precision'' or quantification over a function, the
proof has not yet been compiled to the correct level.

\subsection{Step 8: prove semantic algorithms}

Implement exact reference code and ordinary Lean definitions.  Test malformed
inputs and boundary cases.  These artifacts discover mistakes, but their
status is explicitly ``semantic'', not ``embedded''.

\subsection{Step 9: construct object proofs}

Translate each fixed algorithm to \code{PR}; prove defining equations;
construct local \code{Prf PRAAx} contracts; use the generic trace theorem.
Record the remaining \code{sorry} count and axiom audit.  A file containing
only proposition definitions belongs in the ``targets'' category.

\subsection{Step 10: transport and interpret}

After the \PRAAx\ proof exists, transport it to the native core if desired.
Then prove an external interpretation theorem connecting the codes to the
ordinary theorem.  This last theorem may use a richer Lean mathematics
library without changing what was proved internally.

\subsection{Recommended machine-readable status}

Every packet should contain a ledger with fields:
\begin{lstlisting}[language={}]
{
  "claim": "...",
  "input_presentation": "...",
  "output_relation": "...",
  "producer_status": "proved | conditional | missing | obstructed",
  "verifier_status": "human | semantic_lean | object_prf",
  "precision_transformer": "...",
  "barriers": ["..."],
  "source_commit": "..."
}
\end{lstlisting}

The status values are monotone.  A semantic Lean theorem does not upgrade
itself to \code{object\_prf}; compiled definitions do not upgrade a target to
a theorem; and a finite example does not upgrade a universal claim.

\subsection{Suggested repository library}

A reusable library could be organized as:
\begin{lstlisting}[language={}]
Logic/PrimitiveRecursiveArithmetic/CFP/
  Coding.lean
  BoolFold.lean
  Rational.lean
  FiniteTrace.lean
  Contract.lean
  Demand.lean
  Relational.lean
  BoundedSearch.lean
  FiniteNet.lean
  ObjectCompiler.lean
  Examples/Cubic.lean
\end{lstlisting}
The research packets would import this library and provide only their
domain-specific nodes and producers.  This is the principal engineering gain
of the framework: the difficult proof-theoretic plumbing is done once.
